\par\noindent\begin{minipage}{\linewidth}
\small
\setlength{\tabcolsep}{3pt}
\renewcommand{\arraystretch}{1.15}
\captionof{table}{Persistencia completa tras centrar y con alternativas de dimensión.}\label{tab:S20}
\begin{tabular*}{\linewidth}{@{\extracolsep{\fill}}lrrrrrr@{}}
\toprule
\textbf{Propiedad} & \textbf{Corte} & \textbf{Original} & \textbf{Centrada} & \textbf{\shortstack{Mismos modelos\\y respuestas}} & \textbf{Nuevas} & \textbf{\shortstack{Ahora sin\\resolver}} \\
\midrule
Apertura & 0\% & 262 & 151 & 128 & 9 & 53 \\[2pt]
Apertura & 1\% & 238 & 113 & 78 & 12 & 112 \\[2pt]
Apertura & 5\% & 96 & 11 & 1 & 6 & 91 \\[2pt]
Apertura & 10\% & 19 & 1 & 0 & 1 & 19 \\[2pt]
PR & 0\% & 221 & 211 & 204 & 7 & 16 \\[2pt]
PR & 1\% & 213 & 201 & 194 & 6 & 17 \\[2pt]
PR & 5\% & 177 & 156 & 151 & 4 & 24 \\[2pt]
PR & 10\% & 126 & 106 & 100 & 4 & 23 \\[2pt]
\bottomrule
\end{tabular*}

\par\medskip\textbf{Oposiciones de PR conservadas por las alternativas}\par\smallskip
\begin{tabular*}{\linewidth}{@{\extracolsep{\fill}}lrrrr@{}}
\toprule
\textbf{Corte de PR} & \textbf{PR original} & \textbf{Entropía} & \textbf{D80} & \textbf{Ambas} \\
\midrule
0\% & 221 & 209 & 190 & 190 \\[2pt]
1\% & 213 & 202 & 180 & 180 \\[2pt]
5\% & 177 & 173 & 151 & 151 \\[2pt]
10\% & 126 & 126 & 114 & 114 \\[2pt]
\bottomrule
\end{tabular*}
\par\smallskip{\small Todas las condiciones usan las 26 selecciones y 325 parejas de áreas. Primer bloque: original/centrada son oposiciones totales; mismos modelos y respuestas exige que una pareja concreta original conserve ambas direcciones. Las oposiciones nuevas no lo eran antes. Ahora sin resolver cuenta oposiciones originales que pasan a esa clase. Con corte cero, 142 parejas angulares y 204 de PR conservan la clase de oposición; con los mismos modelos y direcciones son 128 y 204. Ambas direcciones se invierten en 19 parejas de modelos de diez parejas de áreas para apertura, y en ninguna para PR; estos recuentos pueden solaparse con otras categorías. Segundo bloque: el corte se aplica a PR; las alternativas deben conservar dirección no nula en cada selección. Los empates de D80 quedan sin resolver. Los CSV contienen todas las transiciones e identidades.\par}
\end{minipage}\par
